\begin{textsample}{NEG Dim 1 – 2019\_c5 – Score 14.00 – 2019\_c5/enade\_2019\_tec\_seg\_trab\_6283ea4ed7.txt}  \label{ex:f1_neg_391}
CAPÍTULO 5 DISTRIBUIÇÃO DOS CONCEITOS 5.1 PANORAMA NACIONAL DA DISTRIBUIÇÃO DOS CONCEITOS O Conceito Enade é um indicador de qualidade que avalia os cursos por intermédio dos desempenhos dos estudantes a partir dos resultados do Enade.
Os valores possíveis vão de 1 ( pior situação ) a 5 ( melhor situação ).
Os cursos com apenas um ou sem nenhum concluinte participante não obtêm o Conceito Enade, ficando Sem Conceito ( SC ).
Na Tabela 5.1, são apresentadas a quantidade e a distribuição de cursos de Tecnologia em Segurança no Trabalho participantes do Enade/2019, por faixa de conceito e Grande Região.
Enfatiza -se, mais uma vez, que a diferença entre os cursos tabulados neste capítulo e no capítulo 2 corresponde aos cursos que não foram avaliados, ou seja, àqueles, em princípio, sem alunos concluintes que tivessem sido inscritos para a prova.
Observando -se os dados da Tabela 5.1, nota -se que, dos 29 cursos participantes, sete ( 24,1 \% ) classificaram -se com conceito 3, o valor modal.
Esse foi também o conceito modal na região Nordeste, atribuído a cinco cursos ( 41,7 \%, em termos regionais ).
O conceito 2 foi o segundo mais frequente em nível nacional ( 20,7 \%, correspondendo a seis cursos ), e o conceito 5, o terceiro ( 13,8 \%, correspondendo a quatro cursos ).
Houve, ainda, dois cursos ( 6,9 \% ) que receberam conceito 4, e outros dois cursos ( 6,9 \% ) que receberam conceito 1.
Dos 29 cursos de Tecnologia em Segurança no Trabalho, oito ( 27,6 \% ) ficaram Sem Conceito ( SC ).
A região Norte participou com dois cursos ou 6,9 \% do total nacional.
Desses, um curso ( 50,0 \% do total regional ) recebeu conceito 4.
O outro curso avaliado ficou Sem Conceito ( SC ).
Os demais conceitos não foram utilizados.
A região Nordeste participou com 12 cursos ou 41,4 \% do total nacional.
Nessa região, cinco cursos ( 41,7 \% em termos regionais ) obtiveram conceito 3, o conceito modal para a região, como já comentado.
O conceito 5 foi atribuído a três cursos ( 25,0 \% ), os conceitos 1 e 2 foram atribuídos a um curso \textbf{cada} ( 8,3 \% \textbf{cada} ).
Nenhum curso recebeu conceito 4 e dois cursos ( 16,7 \% ) ficou Sem Conceito ( SC ).
Dos doze cursos participantes da região Sudeste ( 41,4 \% do total nacional ), o conceito 2 foi o conceito modal, tendo sido obtido por quatro cursos ( 33,3 \% ).
O conceito 3 foi atribuído a dois cursos ( 16,7 \% ).
Os conceitos 1 e 4 foram atribuídos a um curso, \textbf{cada} ( 8,3 \% \textbf{cada} ), quatro cursos ficaram Sem Conceito ( SC ) e nenhum recebeu o conceito 5.
Os dois cursos da região Sul corresponderam a 6,9 \% do total nacional.
Um curso recebeu conceito 2, já o outro curso avaliado recebeu conceito 5.
Nenhum curso ficou Sem Conceito ( SC ) ou recebeu os demais conceitos.
O único curso participante na região Centro-Oeste ficou Sem Conceito ( SC ).
No Gráfico 5.1, é apresentada a distribuição cumulativa do Conceito Enade segundo Grande Região.
Em um gráfico de distribuição cumulativa, em \textbf{cada} valor de abscissa ( eixo x ), o valor apresentado no eixo das ordenadas ( eixo y ) é o valor acumulado da variável ( conceito Enade ).
Observando a linha vertical pontilhada, que ajuda a visualizar o percentual de conceitos acumulados até o valor 3, ou seja, conceitos 1, 2 e 3, pode -se notar que, por exemplo, a região Sudeste ( linha vermelha ) apresenta 87,5 \% dos cursos com conceito até 3.
A região Sul ( linha rosa ), por outro lado, apresenta 50 \% de seus cursos na mesma situação : com conceito 3 ou abaixo.
A poligonal da região Norte ( linha azul ) mostra que o único curso obteve conceito, este foi igual a 4 e, por isso, a poligonal tem 0 \% no conceito 3 e 100 \% no conceito 4.
Destaca -se que em algumas regiões o conceito 5 não foi alcançado por nenhum curso, como ocorre nas regiões Norte ( linha azul ) e Sudeste ( linha vermelha ), nas quais todos os cursos têm conceito menor ou igual a 4, e as respectivas poligonais cumulativas alcançam 100 \% no conceito 4.
Em linhas gerais, em um gráfico de distribuição cumulativa, poligonais mais à esquerda representam Grandes Regiões com pior distribuição, e poligonais mais à direita, Grandes Regiões com melhores conceitos.
A região Sul ( linha rosa ) apresenta 50 \% de seus cursos com conceito 4 ou abaixo disso.
A região Sudeste apresenta um acúmulo mais suavizado, representando uma distribuição mais diversificada de cursos entre os conceitos, não houve concentração de cursos avaliados por algum conceito específico.
A região Centro-Oeste não é representada no gráfico, seu único curso avaliado ficou Sem Conceito ( SC ). 5.2 CONCEITOS POR CATEGORIA ADMINISTRATIVA, POR MODALIDADE DE ENSINO E GRANDE REGIÃO Na Tabela 5.2, é apresentada a distribuição dos cursos de Tecnologia em Segurança no Trabalho participantes do Enade/2019, por Categoria Administrativa e Modalidade de ensino, de acordo com os conceitos por eles alcançados, segundo as Grandes Regiões brasileiras.
Dos 29 cursos participantes, três ( 10,3 \% ) eram ministrados em Instituições Públicas, e 26 ( 89,7 \% ), em Instituições Privadas.
Quanto à Modalidade de ensino, dez cursos eram a Distância, e 19, Presenciais.
De acordo com as informações da Tabela 5.2, em termos nacionais, entre Instituições Públicas e Privadas, observa -se que, dos quatro cursos avaliados com conceito 5, dois eram oferecidos em IES Públicas, e dois, em IES Privadas.
Dos três cursos participantes de IES Públicas, o conceito 5 foi o valor modal, atribuído aos dois cursos citados ( 66,7 \% ).
O outro curso de IES Públicas foi avaliado com conceito 3.
Nesta categoria, nenhum curso recebeu conceito 1, 2, 4 ou ficou Sem Conceito ( SC ).
Na rede privada, os conceitos modais foram 2 e 3, atribuído a seis cursos \textbf{cada}, o correspondente a 23,1 \% \textbf{cada} conceito, dos 26 cursos da categoria.
Os conceitos 1, 4 e 5 foram atribuídos a dois cursos, \textbf{cada} ( 7,7 \% \textbf{cada} ).
Nesta Categoria Administrativa, oito cursos ( 30,8 \% ) ficaram Sem Conceito ( SC ). foram atribuídos a dois cursos, \textbf{cada} ( 7,7 \% \textbf{cada} ).
Nesta Categoria Administrativa, oito cursos ( 30,8 \% ) ficaram Sem Conceito ( SC ).
Dos 19 cursos participantes da modalidade presencial, o conceito 3 foi o valor modal, atribuído a quatro cursos ( 21,1 \% ).
Dos demais cursos nesta modalidade de ensino, um curso ( 5,3 \% ) recebeu conceito 1, um curso ( 5,3 \% ) recebeu conceito 2, dois cursos ( 10,5 \% ), conceito 4, e três cursos ( 15,8 \% ), conceito 5.
Nesta modalidade de ensino, oito cursos ( 42,1 \% ) ficaram Sem Conceito ( SC ).
Os dez cursos a Distância informados foram avaliados nos conceitos 1, 2, 3 e 5.
Os cursos ficaram distribuídos da seguinte maneira : conceito 1 ( um curso ), 2 ( cinco cursos ), 3 ( três cursos ) e 5 ( um curso ).
Nenhum curso ficou Sem Conceito ( SC ) ou recebeu conceito 4.
Na análise por região, observa -se que, na região Norte, as Instituições Privadas participaram com os dois cursos da região ( 100,0 \% do total regional ), dos quais um ficou Sem Conceito ( SC ), e um recebeu conceito 4.
Nenhum recebeu os demais conceitos.
As Instituições Públicas não tiveram cursos avaliados.
Todos os cursos oferecidos na região Norte eram presenciais.
Na região Nordeste, a rede privada concentrou nove dos 12 cursos participantes, o equivalente a 75,0 \% do total da região.
Dentre os cursos oferecidos por IES Privadas no Nordeste, quatro cursos foram avaliados com conceito 3, o conceito modal.
Os conceitos 1, 2 e 5 foram atribuídos a um curso, \textbf{cada}.
Nenhum curso recebeu conceito 4 e dois cursos ficaram Sem Conceito ( SC ).
As Instituições Públicas da região Nordeste participaram com três cursos ( 25,0 \% do total da região ).
Desses, dois receberam conceito 5, o valor modal, seguidos por um curso, que recebeu conceito 3.
Nenhum curso ficou Sem Conceito ( SC ) ou com os demais conceitos nessa região.
Nove cursos eram presenciais na região Nordeste, esses receberam os conceitos 2 ( um curso ), 3 e 5 ( três cursos, \textbf{cada} ).
Os cursos a distância eram três, dois receberam conceito 3 e um curso recebeu conceito 1.
Na região Sudeste, todos os oito oferecidos eram de rede privada ( 100,0 \% em termos regionais ).
Quatro cursos receberam conceito 2, dois cursos receberam conceito 3, os conceitos 1 e 4 foram atribuídos a um curso, \textbf{cada}.
Quatro cursos ficaram Sem Conceito ( SC ) e nenhum recebeu conceito 5.
As IES Públicas não tiveram cursos avaliados.
Também na região Sudeste, a maioria dos cursos foi oferecida na modalidade presencial : sete cursos.
Para esses, três cursos obtiveram os conceitos 1, 3 e 4.
Quatro cursos ficaram Sem Conceito ( SC ) e nenhum recebeu os demais conceitos.
Na região Sudeste, os cinco cursos oferecidos na modalidade Educação a Distância receberam os conceitos 2 ( quatro cursos ) e 3 ( 1 curso ).
Nenhum curso ficou Sem Conceito ( SC ) ou recebeu os demais conceitos.
Os dois cursos avaliados na região Sul faziam parte de IES Privadas.
Desses, um recebeu conceito 2 e o outro, conceito 5.
As IES Públicas não tiveram cursos avaliados, como dito anteriormente.
Ambos os cursos oferecidos na região Sul eram a distância.
Na região Centro-Oeste, apenas um curso foi avaliado, e pertencia a uma IES Privada.
Tal curso era presencial e ficou Sem Conceito ( SC ).
No Gráfico 5.2, apresenta -se a distribuição do Conceito Enade, segundo a Categoria Administrativa da IES.
Os cursos em IES Públicas ( linha azul ) apresentam uma poligonal mais à direita do que os em IES Privadas ( linha verde ) e, consequentemente, uma distribuição de Conceitos Enade com valores maiores. 5.3 CONCEITOS POR ORGANIZAÇÃO ACADÊMICA E GRANDE REGIÃO Na Tabela 5.3, encontra -se a distribuição dos conceitos atribuídos aos cursos participantes do Enade/2019, na Área de Tecnologia em Segurança no Trabalho, por Organização Acadêmica, segundo as Grandes Regiões brasileiras.
Dos 29 cursos de Tecnologia em Segurança no Trabalho participantes, onze eram oferecidos em Universidades, dez em Centros Universitários, sete em Faculdades e um em CEFET/IF.
Esta distribuição corresponde a, respectivamente, 37,9 \%, 34,5 \%, 24,1 \% e 3,4 \% dos cursos.
De acordo com os dados apresentados, de todos os quatro cursos avaliados com conceito 5, um era vinculado a Universidades.
Os cursos vinculados a esse tipo de Organização Acadêmica tiveram o conceito 2 como conceito modal, atribuído a quatro cursos ( 36,4 \% ).
Os demais cursos vinculados a Universidades receberam os conceitos 1 ( dois cursos ), 3 ( três cursos ) e conceito 5 ( um curso ).
Um curso ficou Sem Conceito ( SC ) e nenhum recebeu conceito 4.
Dos cursos em Centros Universitários, o conceito modal foi 3, atribuído a três cursos ( 30,0 \% ).
Dois cursos vinculados a esse tipo de Organização Acadêmica receberam conceito 4.
Os conceitos 2 e 5 foram atribuídos a um curso, \textbf{cada}.
Três cursos ficaram Sem Conceito e nenhum recebeu conceito 1.
Dos sete cursos mantidos por Faculdades, os conceitos 2, 3 e 5 foram atribuídos a um curso, \textbf{cada} ( 14,3 \% \textbf{cada} ).
Quatro cursos ficaram Sem Conceito ( SC ) e nenhum recebeu conceitos 1 e 4.
O único curso mantido por CEFET/IF recebeu conceito 5.
Considerando -se, separadamente, as regiões brasileiras, verifica -se que, na região Norte, as Universidades e CEFET/IF não tiveram cursos avaliados.
Os Centros Universitários da região Norte foram representados por apenas um curso, o qual obteve conceito 4.
As Faculdades também participaram com apenas um curso, que ficou Sem Conceito ( SC ).
Na região Nordeste, as Universidades participaram com três dos 12 cursos da Área de Tecnologia em Segurança no Trabalho oferecidos.
Os conceitos 1, 3 e 5 foram atribuídos a um curso, \textbf{cada}.
Nenhum curso ficou Sem Conceito ( SC ) ou recebeu os demais conceitos.
Os Centros Universitários contaram com quatro cursos participantes na região Nordeste, três receberam conceito 3 e um ficou Sem Conceito ( SC ).
As Faculdades foram representadas por quatro cursos na região Nordeste.
Os conceitos 2, 3 e 5 foram atribuídos a um curso, \textbf{cada}.
Um curso ficou Sem Conceito ( SC ).
Nenhum curso recebeu os conceitos 1 e 4.
Os CEFET/IF participaram com apenas um curso, que curso recebeu conceito 5.
Na região Sudeste, as Universidades concentraram oito dos doze cursos de Tecnologia em Segurança no Trabalho da região.
Dos cursos oferecidos em Universidades, nessa região, o conceito modal foi 2, atribuído a quatro cursos.
Os demais cursos receberam os conceitos 1 ( um curso ) e 3 ( dois cursos ).
Um curso ficou Sem Conceito ( SC ) e nenhum foi avaliado com os conceitos 4 e 5.
Os Centros Universitários participaram com três cursos.
Um recebeu o conceito 4 e dois ficaram Sem Conceito ( SC ).
As Faculdades participaram com apenas um curso e os CEFET/IF não tiveram cursos avaliados.
Os dois cursos da região Sul eram mantidos por Centros Universitários.
Um curso recebeu conceito 2, e o outro, conceito 5.
Universidades, Faculdades e CEFET/IF não tiveram cursos avaliados.
A região Centro-Oeste teve apenas um curso avaliado, que era mantido por Faculdades, porém tal curso ficou Sem Conceito ( SC ).
Universidades, Centros Universitários e CEFET/IF não tiveram cursos avaliados.
No Gráfico 5.3, apresenta -se a distribuição do Conceito Enade dos cursos de Tecnologia em Segurança no Trabalho, segundo a Organização Acadêmica da IES.
Os cursos em CEFET/IF ( linha laranja ) apresentam uma poligfoional mais à direita do que as demais, denotando uma melhor distribuição de Conceitos Enade.
No outro extremo, com distribuição muito semelhante entre si, encontram -se os cursos mantidos por Universidades ( linha azul ) e Faculdades ( linha vermelha ), apresentando as respectivas poligonais mais à esquerda e, consequentemente, distribuição de Conceitos Enade com menores valores.

% matched lemmas: cada
\end{textsample}
